%% ══════════════════════════════════════════════════════════
%%  ITEM — Shuffled Time   Sponsorship Proposal  (English)
\newcommand{\sponsor}{BALCIA}
%% ══════════════════════════════════════════════════════════

\documentclass[a4paper,9.5pt]{article}

\usepackage[top=0cm,bottom=0cm,left=0cm,right=0cm]{geometry}
\usepackage{xcolor}
\usepackage[T1]{fontenc}
\usepackage[utf8]{inputenc}
\usepackage{tikz}
\usepackage{graphicx}
\usepackage{tabularx}
\usepackage{parskip}
\usepackage{microtype}
\usepackage{lmodern}
\usepackage{textcomp}
\usepackage{enumitem}
\usepackage{multicol}
\usepackage{ifthen}
\usepackage{hyperref}

\hypersetup{hidelinks}
\pagestyle{empty}
\setlength{\parskip}{1.5pt}
\setlength{\columnsep}{0.8em}

\definecolor{raacream}{HTML}{F5F5DC}
\definecolor{raared}{HTML}{CC2222}
\definecolor{raadark}{HTML}{0D0D0D}
\definecolor{itemgold}{HTML}{D4A84B}
\definecolor{mygray}{HTML}{555555}

\newcommand{\sectionrule}[1]{%
  \vspace{3pt}%
  {\color{raared}\bfseries\normalsize #1}\\[-3pt]%
  {\color{raared}\rule{\linewidth}{0.8pt}}%
  \vspace{2pt}%
}

%% Performance card: {date}{title}{artist}{img-or-none}{description}{bio}
\newcommand{\perfcard}[6]{%
  \noindent\begin{minipage}[t]{0.37\linewidth}
    \vspace{0pt}
    {\small\color{raared}\textbf{#1}}\\[1pt]
    {\bfseries\small #2}\\[2pt]
    {\footnotesize\textit{#3}}\\[5pt]
    \ifthenelse{\equal{#4}{none}}{%
      {\setlength{\fboxsep}{0pt}%
       \colorbox{black!8}{\parbox[c][3.6cm][c]{\linewidth}{%
         \centering\footnotesize\color{mygray}impression to follow}}}%
    }{%
      \makebox[\linewidth][c]{%
        \includegraphics[height=3.6cm,width=\linewidth,keepaspectratio=true]{#4}}%
    }%
  \end{minipage}%
  \hfill%
  \begin{minipage}[t]{0.60\linewidth}
    \vspace{0pt}
    {\small #5}%
    \par\vspace{5pt}%
    {\footnotesize\color{mygray}\textbf{Bio:}\enspace #6}%
  \end{minipage}%
  \par\vspace{0.35cm}%
  {\color{raared!35}\rule{\linewidth}{0.4pt}}%
  \par\vspace{0.35cm}%
}

\begin{document}

%% ══════════════════════════════════════════════════════════
%%  PAGE 1 — FULL-BLEED POSTER + PROPOSAL CONTENT
%% ══════════════════════════════════════════════════════════
\noindent
\begin{tikzpicture}[x=1cm,y=1cm]
  \useasboundingbox (0,0) rectangle (21,29.7);

  %% Poster stretched edge-to-edge, no margins
  \node[anchor=south west, inner sep=0pt] at (0,0) {%
    \includegraphics[width=21cm,height=29.7cm,keepaspectratio=false]{img/item-poster.jpg}%
  };

  %% Dark overlay for text legibility
  \fill[black,opacity=0.67] (0,0) rectangle (21,29.7);

  %% ── HEADER ──────────────────────────────────────────────

  \node[anchor=north west, inner sep=0pt] at (0.9,29.4) {%
    
\includegraphics[height=1.1cm]{img/raa_logo_v2.png}%
  };

  \node[anchor=north east, inner sep=0pt] at (20.1,29.5) {%
    \fontsize{5.5}{7}\selectfont\color{raacream!50}%
    \textls[180]{S P O N S O R S H I P \ \ P R O P O S A L}%
  };

  \node[anchor=center] at (10.5,26.8) {%
    \fontsize{26}{31}\selectfont\bfseries\color{raacream}%
    ITEM\,—\,SHUFFLED TIME%
  };

  \node[anchor=center] at (10.5,25.9) {%
    \fontsize{9}{11}\selectfont\color{raacream!80}%
    \textit{5.–11. July 2026\enspace·\enspace RAA.SPACE, Matīsa iela 8, Rīga}%
  };

  \node[anchor=center] at (10.5,24.85) {%
    \fontsize{9}{11}\selectfont\color{itemgold}%
    Presented to\enspace\textbf{\sponsor}%
  };

  \fill[raared] (0,24.15) rectangle (21,24.45);

  %% ── BODY ─────────────────────────────────────────────────
  \node[anchor=north west, text width=19cm, inner sep=0pt] at (1.0,24.0) {%
    \begin{minipage}[t]{19cm}
      \color{raacream}
      \setlength{\parskip}{1.5pt}
      \setlength{\columnsep}{0.8em}

      \vspace{0.18cm}
      \sectionrule{About the Event}

      \small
      ITEM\,—\,Shuffled Time is a seven-day performative week at RAA.SPACE in Rīga.
      Each day one artist or artist group occupies RAA's main hall — a 138\,m\textsuperscript{2}
      space with ten windows facing Matīsa iela — for a 24-hour period to \textbf{I}nvoke,
      \textbf{T}emper, \textbf{E}late, \textbf{M}int the space. Performances may span the full
      24 hours or take place over a shorter duration with one or more openings. The programme
      brings together seven distinct artistic voices across durational performance, body art,
      experimental poetry, aromagraphy and interactive installation.
      Organised by RAA and supported by Veldze Radošā Rūpnīca.

      \vspace{0.1cm}

      \begin{multicols}{2}

        \sectionrule{Programme}

        {\footnotesize\setlength{\parskip}{0pt}
        \begin{tabular}{@{}p{1.35cm}p{5.65cm}@{}}
        \textbf{Sun 5.07}   & \textbf{Ilze Aulmane}\\
                             & \textit{(details to follow)}\\[1pt]
        \textbf{Mon 6.07}   & \textbf{NOW} — Maruta Gredzena\\
                             & Interactive wearable performance\\[1pt]
        \textbf{Tue 7.07}   & \textbf{N.3} — Ernests Ā\\
                             & Rope \& endurance, 18:00–24:00\\[1pt]
        \textbf{Wed 8.07}   & \textbf{Gluttony} — K.\,Šuksta \& S.\,Murins\\
                             & Aromagraphy performance\\[1pt]
        \textbf{Thu 9.07}   & \textbf{Presence in Absence Pt.\,3}\\
                             & Lisette Ros\\[1pt]
        \textbf{Fri 10.07}  & \textbf{Time Scores} — Vanessa Brzič\\
                             & Fluxus-score participatory piece\\[1pt]
        \textbf{Sat 11.07}  & \textbf{Laiks Telpā} — Ilze Mazpane\\
                             & \& Jana Zariņa\\
                             & Durational music-body-space performance\\
        \end{tabular}
        }

        \columnbreak

        \sectionrule{Audience \& Reach}

        {\small
        \begin{itemize}[leftmargin=*,itemsep=1pt,topsep=0pt]
          \item \textbf{200+} visitors across the week
          \item International and local artists, art professionals, curious public
          \item Instagram \textbf{@raa\_riga} \& \textbf{@veldze.radosarupnica}:\\
                combined \textbf{7\,000+} followers
          \item Media reach: \textbf{100gkultura}, \textbf{riga-neighbourhood},\\
                \textbf{arterritory.com}, \textbf{notboring.riga}
        \end{itemize}
        }

        \vspace{0.05cm}
        \sectionrule{Why It Matters}

        {\small Live and durational performance is raw and unrepeatable.
        ITEM places \textbf{\sponsor} within RAA.SPACE — a space with an established
        programme of performance and visual art since its founding, with collaborations
        including Homo Novus, Kefīrs Studio, 2 Annas, Alīna Burlakova, Andrejs Osokins
        and many more. Sponsorship here is a genuine cultural investment, seen and felt
        by an engaged, arts-literate audience.}

      \end{multicols}

      \vspace{0.05cm}

      \sectionrule{Sponsorship Package — \texteuro\,2\,000}

      {\small
      \begin{tabularx}{19cm}{@{}>{\bfseries}p{3.3cm}X@{}}
      Credited everywhere   & \sponsor{} acknowledged at every event, on all print materials
                              and on the RAA.SPACE event page.\\[2pt]
      Social media          & Dedicated posts \& stories across @raa\_riga,
                              @veldze.radosarupnica, riga-neighbourhood, notboring.riga,
                              100gkultura, arterritory.com (7\,000+ combined reach).\\[2pt]
      VIP access            & 3 reserved seats per evening. Invitation to artist
                              meet-and-greets and opening nights.\\[2pt]
      Speaking opportunity  & Brief welcome address at the opening evening (5 Jul).\\[2pt]
      On-site presence      & \sponsor{} display in the RAA foyer throughout the week.\\[2pt]
      Press coverage        & Named as principal sponsor in all media communications.\\
      \end{tabularx}
      }

    \end{minipage}%
  };

  %% ── FOOTER STRIP ─────────────────────────────────────────
  \fill[raadark,opacity=0.88] (0,0) rectangle (21,0.82);
  \node[anchor=center] at (10.5,0.41) {%
    \fontsize{7.5}{9}\selectfont\color{raacream!80}%
    \textbf{RAA.SPACE}\enspace·\enspace
    Matīsa iela 8, Rīga\enspace·\enspace
    \href{https://raa.space}{raa.space}\enspace·\enspace
    \href{mailto:enter@raa.space}{enter@raa.space}\enspace·\enspace
    @raa\_riga%
  };

\end{tikzpicture}

%% ══════════════════════════════════════════════════════════
%%  PAGES 2+ — PERFORMANCE CARDS
%% ══════════════════════════════════════════════════════════
\newpage
\newgeometry{top=1.5cm,bottom=1.5cm,left=1.5cm,right=1.5cm}

\noindent
{\color{raared}\bfseries\large ITEM\,—\,SHUFFLED TIME}%
{\small\color{mygray}\quad Performance Programme · 5–11 July 2026 · RAA.SPACE, Rīga}%
\\[-4pt]%
{\color{raared}\rule{\linewidth}{0.8pt}}

\vspace{0.45cm}

%% ── 1 · Ilze Aulmane ─────────────────────────────────────
\perfcard%
  {Sunday, 5 July}%
  {Ilze Aulmane}%
  {Ilze Aulmane}%
  {img/item-aulmane.jpg}%
  {Full title, description and visual impression will be shared shortly.}%
  {Bio to follow.}%

%% ── 2 · NOW ──────────────────────────────────────────────
\perfcard%
  {Monday, 6 July}%
  {NOW}%
  {Maruta Gredzena}%
  {img/item-now.jpg}%
  {In her interactive performance \textit{Now}, Maruta Gredzena explores whether
   communication without words can be true. Visitors become participants, choosing
   and wearing light, paper-based wearable sculptures in response to their present
   state — soft, fluffy, geometric, or rock-like forms creating a temporary second
   skin. Each choice reshapes a shifting collective composition where everyone
   becomes both observer and artwork. The work asks: can honesty exist without
   language? A paper sculpture installation with objects and strips descending from
   the ceiling provides the spatial backdrop.}%
  {Maruta Gredzena is interested in what happens when nothing stays still. A life
   shaped by continual relocation has shifted her practice from making objects to
   creating spaces for interaction. Rooted in porcelain and ceramic sculpture and
   expanded through graphic art, her work unfolds through installations that use
   paper, material, and playfulness to explore communication beyond language.
   IG:~@em.gaa}%

%% ── 3 · N.3 ──────────────────────────────────────────────
\perfcard%
  {Tuesday, 7 July}%
  {N.3}%
  {Ernests Ā}%
  {img/item-n3.jpg}%
  {\textit{The binding conviction in the existence of two opposites always makes
   one wander in confusion.}\par\vspace{3pt}%
   Ernests arrives with a rope, initially wrapping it around one of the gallery's
   columns. He ties the end to his leg and walks around the column in widening
   circles — both within and beyond the gallery walls — attempting to distance
   himself from constraint until physical exhaustion or external factors intervene.
   Should he break free, he may wrap the rope once again. Brief rests do not
   interrupt the piece. He remains silent throughout. Intended duration:
   18:00–24:00.}%
  {In his practice, Ernests Ā devotes himself to playful exploration of the mundane
   and undefinable aspects of reality within unrestricted forms of expression.}%

%% ── 4 · Gluttony ─────────────────────────────────────────
\perfcard%
  {Wednesday, 8 July}%
  {Gluttony}%
  {Kristiāna Šuksta \& Sandris Murins}%
  {img/item-gluttony.jpg}%
  {An experimental aromagraphy performance pairing twelve words — each coupled with
   a distinct aroma — performed in continuously reshuffled sequences and tempos.
   Through congruent and incongruent combinations of language and scent,
   \textit{Gluttony} reveals how olfactory experience can reinforce, contradict or
   transform the meaning of text — exploring desire, memory and the sensory
   dimensions of meat consumption. Meaning emerges through ever-changing temporal
   configurations rather than a fixed sequence.}%
  {Sandris Murins is a Latvian artist, inventor of the Scent Camera, and founder of
   the Scent + Art Festival — promoting aromagraphy as an emerging artistic language.
   Kristiāna Šuksta is a Latvian poet, performance artist and technical producer of
   the festival Starptelpa. Her practice combines poetry and live performance to
   explore personal experience and human relationships through wit and humour.}%

%% ── 5 · Presence in Absence, Piece Three ─────────────────
\perfcard%
  {Thursday, 9 July}%
  {Presence in Absence, Piece Three}%
  {Lisette Ros (Dutch)}%
  {img/item-presence.jpg}%
  {\textit{Presence in Absence} is a series by Lisette Ros begun in Mexico City
   in 2023, examining the perceptible act of presence in absence through
   \textit{The Absent Body}. In Piece Three, Ros grapples with her own presence
   by taking in a natural hypnotic — a substance that has formed part of her
   nightly ritual since 2014, when she was diagnosed with DSPS (Delayed Sleep
   Phase Syndrome). The work manipulates control and consciousness in real time,
   before the audience.}%
  {Lisette Ros (b.\,1991, NL) is a Dutch conceptual and performance artist with
   Indonesian roots. Her work investigates identity as a construct formed through
   social systems and embodied experience, using the body as research tool and
   medium. Presented internationally at the Venice Architecture Biennale, SXSW,
   Van Abbemuseum, Goethe Institute, MWoods Beijing, Casa Lü CDMX. 2021 residency
   with Marina Abramović.}%

%% ── 6 · Time Scores ──────────────────────────────────────
\perfcard%
  {Friday, 10 July}%
  {Time Scores}%
  {Vanessa Brzič (Croatian)}%
  {img/item-timescores.jpg}%
  {Audience members draw cards with short instructions — such as ``cut'', ``tie'',
   ``spill'', ``listen'' — and carry them out in the space using various objects.
   Each action changes the room for the next one, so the same card means something
   different every time it is drawn. The room becomes a record of everything that
   has happened inside it.}%
  {Vanessa Brzič is a cross-disciplinary artist and OIST QISS Fellow working at
   the intersection of fundamental research, quantum information and artistic
   practice. Details being confirmed.}%

%% ── 7 · Laiks Telpā ─────────────────────────────────────
\perfcard%
  {Saturday, 11 July}%
  {Laiks Telpā}%
  {Ilze Mazpane \& Jana Zariņa}%
  {img/item-laikstelpa.jpg}%
  {A durational performance that, using only objects found on site within RAA's
   existing space, creates a musical-bodily-spatial experience of subjective time.
   The work explores how time expands, contracts, rushes, stops and disappears —
   without clock-time as guide or constraint.}%
  {Ilze Mazpane is a performance artist from Rīga with a background in philosophy,
   visual art, experimental theatre and creative writing. She creates site-specific
   durational works marked by poetic sensibility and a sensitive relationship with space.
   Jana Zariņa is primarily a historical keyboard and sound artist whose work spans
   early music, improvisation, electro-acoustic performance and interdisciplinary arts.
   Key projects include the electro-acoustic duo \textit{maiak}, the contemporary
   Latvian opera \textit{Ticks} (2024, Rīga Zoo) and international baroque
   performances. She has performed with the Latvian National Symphony Orchestra.}%

\end{document}
